\scanpage{54}
\S\ 4 \uline{Topos stratifiant et topos cristallin.}

4.1. Interprétation d'une stratification relative d'un Module $\underline{F}$ sur $X/S$.
Si \unsure{on a} une telle stratification, alors pour tt morphisme $f:\unsure{Y}\to X$, et tt voisinage infinitésimal $Y'$ de $Y$, \inserted{(sur $S$)} \deleted{\unsure{qui que}} tel que $f$ veut bien se prolonger (localement au moins) en $f':Y'\mathbin{-}X$, un système \uline{transitif} d'isom entre les $f'^*(\underline{F})$.
\margin{\diagram{\begin{tikzcd}[ampersand replacement=\&]
X \arrow[dd,dash] \& \unsure{Y_0|Y} \arrow[l,"f"'] \arrow[d,{Hooks[right]}-] \\
\& \unsure{Y'_0|Y'} \arrow[ul,"f'"]\\
S \&
\end{tikzcd}}}
On peut donc dire que la \deleted{\unsure{donnée}} \inserted{donnée d'une} stratification \deleted{\unsure{nous}} \inserted{revient à la} donnée \deleted{\unsure{canonique}}, pour tt voisinage infinitésimal $Y'$ d'un $Y\xrightarrow{f}X$, (sur $S$) \inserted{soumis à la condition déjà énoncée (autom. si $X/S$ form$^{\mathrm{t}}$ lisse)} d'un \uline{prolongement} \deleted{\unsure{can.}} de $F$ en un \uline{Module} sur $Y'$, avec des conditions de fonctorialité évidentes.
Pour exprimer plus commodément ce point de vue, et ses variantes, on va introduire des sites ad-hoc.

\margin{53}
